% Section 25.6: Electromagnetism: which fields are $1$-forms, $2$-forms, $3$-forms? | Part: 5 | Depends on: part:manifolds, part:integration, sec:lc | Status: todo | Original lines: 3559-3608
%!TEX root = ../../../main.tex
\subsection{Electromagnetism: which fields are $1$-forms, $2$-forms, $3$-forms?}\label{sec:phys-em}

With $\R^3$ as space, the natural degrees are:
\begin{center}
\small
\renewcommand{\arraystretch}{1.3}
\begin{tabular}{@{}p{5.2cm}p{1.8cm}p{6.0cm}@{}}
\toprule
\textbf{Field} & \textbf{Degree} & \textbf{Measured by}\\
\midrule
$E$ (electric field) & $1$-form & $\oint_\gamma E$ (e.m.f.), work per unit charge\\
$H$ (magnetic field strength) & $1$-form & $\oint_\gamma H$ (magnetomotive force)\\
$B$ (magnetic flux density) & $2$-form & $\int_SB$ (magnetic flux)\\
$D$ (electric displacement) & $2$-form & $\int_SD$ (electric flux)\\
$\rho$ (charge density) & $3$-form & $\int_\Omega\rho$ (charge)\\
$j$ (current density) & $2$-form & $\int_Sj$ (current)\\
$A$ (potential) & $1$-form & $\oint_\gamma A$ (holonomy / phase)\\
\bottomrule
\end{tabular}
\end{center}
($D$, $H$, $\rho$, $j$ are in fact \emph{twisted} forms: they also need an outer orientation, see Hehl--Obukhov or Burke.) The integral laws are then metric-free, being pure applications of Stokes' theorem:
\begin{alignat*}{2}
\oint_{\partial S}E&=-\frac{\dd}{\dd t}\int_SB&&\iff\dd E+\partial_tB=0\quad\text{(Faraday),}\\
\int_{\partial\Omega}B&=0&&\iff\dd B=0\quad\text{(no monopoles),}\\
\int_{\partial\Omega}D&=\int_\Omega\rho&&\iff\dd D=\rho\quad\text{(Gauss),}\\
\oint_{\partial S}H&=\int_Sj+\frac{\dd}{\dd t}\int_SD&&\iff\dd H=j+\partial_tD\quad\text{(Amp\`ere--Maxwell).}
\end{alignat*}
\emph{From integral to differential form (Faraday as an example).} By Stokes, $\oint_{\partial S}E=\int_S\dd E$; for a surface $S$ fixed in time, $-\frac{\dd}{\dd t}\int_SB=-\int_S\partial_tB$. So the law says $\int_S(\dd E+\partial_tB)=0$ for every surface $S$; a $2$-form whose integral over every small surface patch vanishes is itself zero, hence $\dd E+\partial_tB=0$. The same argument gives the other three.

The metric enters only through the \emph{constitutive relations} of the medium (vacuum): $D=\varepsilon_0\hodge E$, $B=\mu_0\hodge H$. Hence Gauss' law is topological, and the factor $E\cdot n\,\dd A$ of the usual statement hides $\varepsilon_0\hodge E$. In spacetime, $F=E\wedge\dd t+B$ and $\dd F=0$ is metric-free, while $\dd\hodge F=J$ needs $\hodge$ (in dimension $4$ the star on $2$-forms depends only on the conformal class of the metric, so Maxwell's equations are conformally invariant). Charge conservation is $\dd J=0$, and by Stokes $\int_{\partial\Omega}J=0$.
(In geometric algebra the two spacetime equations are the single equation $\vd F=J$, Section~\ref{sec:synth2}.)

\subhead{The Lorentz force and the interior product}
The force on a charge $q$ with velocity $v$ is the covector
\begin{equation}\label{eq:lorentz}
F=q\bigl(E-\iota_vB\bigr),\qquad\text{i.e. }F(w)=q\bigl(E(w)-B(v,w)\bigr).
\end{equation}
\emph{Derivation in $\R^3$ (steps).} Let $B=\iota_{\mathbf B}\vol$, \ie $B(a,b)=\vol(\mathbf B,a,b)$. The magnetic part of the Lorentz force is $F_{\mathrm{mag}}(w)=q\,(v\times\mathbf B)\cdot w$.
\begin{enumerate}
\item The scalar triple product is the volume form: $(v\times\mathbf B)\cdot w=\vol(v,\mathbf B,w)$.
\item Swap the first two arguments (alternating): $\vol(v,\mathbf B,w)=-\vol(\mathbf B,v,w)=-B(v,w)$.
\item Hence $F_{\mathrm{mag}}(w)=-q\,B(v,w)=-q\,(\iota_vB)(w)$, which proves the formula.
\end{enumerate}
The power is $F(v)=qE(v)-qB(v,v)=qE(v)$, because $B$ is alternating ($B(v,v)=0$): the magnetic force does no work for the same reason that $\iota_v\iota_v=0$.

\begin{example}[Amp\`ere's law and potentials: the angle form again]\label{ex:ampere}
Outside a wire carrying current $I$, the field has magnitude $I/(2\pi r)$ and points along $\hat\theta$, so, as a covector, $H=\frac I{2\pi r}e^\theta=\frac I{2\pi r}r\,\dd\theta=\frac I{2\pi}\dd\theta$ (using $e^\theta=r\,\dd\theta$). This is closed but not exact: $\oint H=I\neq0$, and there is no single-valued scalar potential even though $\dd H=0$.
The Aharonov--Bohm effect is the same phenomenon in quantum mechanics: the phase $\frac q\hbar\oint A$ along a loop around a solenoid is non-zero although $B=\dd A=0$ at every point of the loop. In the physical setup Stokes holds: $\oint A=\int_SB=\Phi$, the flux through the solenoid, because every surface $S$ bounded by the loop must cross the solenoid, where $B\neq0$. In the idealised thin-solenoid model the solenoid is removed from space; then $A$ is closed but not exact on the remaining region, exactly like $\frac I{2\pi}\dd\theta$ above, and no surface inside that region is bounded by the loop.
\end{example}
